\chapter{Tuberculosis TB}
\index{Tuberculosis TB}

\section*{Definition and Overview}
Tuberculosis (TB) is a chronic, communicable infectious disease caused by pathogens of the \textit{Mycobacterium tuberculosis} complex. It represents one of the oldest recorded diseases of humanity, historically referred to as ``consumption,'' the ``white plague,'' or ``phthisis'' due to the profound muscle wasting and cachexia seen in advanced cases. Despite the availability of effective, curative antibiotic regimens for over seven decades, tuberculosis remains a major global public health crisis and one of the leading infectious causes of death worldwide, carrying significant morbidity and mortality.

The clinical manifestations of tuberculosis are classically categorized into pulmonary and extrapulmonary forms. Pulmonary tuberculosis, which affects the lung parenchyma, is the most common presentation and the sole source of transmission. Extrapulmonary tuberculosis occurs when the bacilli disseminate via the bloodstream or lymphatic system to other organ systems, including the lymph nodes, pleura, central nervous system, bones and joints, genitourinary tract, and gastrointestinal tract. Extrapulmonary TB is non-infectious unless concurrent pulmonary involvement is present or a procedure aerosolizes the pathogen.

A central paradigm of tuberculosis is the distinction between Latent Tuberculosis Infection (LTBI) and Active Tuberculosis Disease. Latent TB is characterized by a state of persistent immune response to mycobacterial antigens without evidence of clinically manifest active disease. Individuals with LTBI are asymptomatic, have normal chest radiographs, and are non-infectious. However, they harbor viable, dormant bacilli within microscopic granulomas, representing a reservoir for potential reactivation. Active TB disease, on the other hand, is characterized by the uncontrolled replication and spread of the bacilli, resulting in tissue destruction, clinical symptoms, and, in pulmonary cases, infectivity to others. Controlling the transition from latent infection to active disease is a primary focus of modern tuberculosis eradication strategies.

\section*{Epidemiology}
Tuberculosis is a global epidemic that exhibits a highly disproportionate burden in low- and middle-income nations. According to estimates by the World Health Organization (WHO), approximately one-quarter of the global population is infected with latent \textit{Mycobacterium tuberculosis}, representing a massive pool of potential future active cases. Annually, an estimated 10 million to 10.6 million individuals develop active tuberculosis disease, resulting in approximately 1.3 million to 1.5 million deaths. This includes approximately 167,000 deaths among people living with HIV (PLHIV), making TB the leading cause of death among this immunocompromised cohort.

Geographically, the burden is concentrated in specific regions. The WHO identifies 30 high-TB-burden countries that collectively account for over 86\% of all global cases. The South-East Asia region bears the largest proportion of the global burden (approximately 45\%), followed by the African region (23\%) and the Western Pacific region (18\%). Individual countries with the highest absolute number of active TB cases include India, Indonesia, China, the Philippines, Pakistan, Nigeria, Bangladesh, and the Democratic Republic of the Congo.

Within both endemic and non-endemic countries, specific populations are disproportionately affected by tuberculosis due to biological, environmental, and socioeconomic factors:
\begin{itemize}
    \item \textbf{Immunocompromised Cohorts:} Individuals with impaired cell-mediated immunity, particularly those living with HIV, are at a 15- to 20-fold higher risk of progressing from latent to active TB. Other immunocompromised groups include patients undergoing tumor necrosis factor (TNF) inhibitor therapy, chemotherapy, or organ transplant immunosuppression.
    \item \textbf{Socioeconomically Disadvantaged Populations:} Poverty, malnutrition, homelessness, and overcrowding significantly increase the risk of transmission and progression. Overcrowded environments with poor ventilation, such as correctional facilities, refugee camps, and informal settlements, serve as hotbeds for TB transmission.
    \item \textbf{Migrants and Refugees:} Populations migrating from high-burden regions to low-incidence countries often display higher rates of latent and active TB, reflecting the epidemiology of their regions of origin rather than their host countries.
    \item \textbf{Age and Gender Distributions:} Active TB is most prevalent among young adults in their economically productive years (ages 15 to 59). However, infants, young children, and the elderly are at a higher risk of developing severe, disseminated forms of the disease. Globally, males account for approximately 56\% of all notified TB cases, showing a male-to-female ratio of 1.6:1, which is attributed to both biological susceptibility and gender-based exposure differences.
\end{itemize}

The rise of drug-resistant tuberculosis is a critical challenge. Multi-drug resistant TB (MDR-TB) is defined as resistance to at least isoniazid and rifampicin, the two most potent first-line drugs. Extensively drug-resistant TB (XDR-TB) is defined as MDR-TB with additional resistance to any fluoroquinolone and at least one additional Group A drug (such as bedaquiline or linezolid). Approximately 400,000 cases of MDR-TB or rifampicin-resistant TB occur annually, posing a severe threat to global health due to the complexity, cost, and toxicity of second-line therapies.

\section*{Aetiology and Pathophysiology}
Tuberculosis is caused by infection with \textit{Mycobacterium tuberculosis}, a slow-growing, obligate aerobic, non-motile, non-spore-forming, acid-fast bacillus. The bacterium belongs to the \textit{Mycobacterium tuberculosis} complex, which also includes \textit{M. bovis}, \textit{M. africanum}, \textit{M. microti}, and \textit{M. canettii}. The distinguishing feature of the mycobacterial cell wall is its high lipid content, particularly mycolic acids. This complex, waxy outer layer protects the bacterium from mechanical and chemical injury, prevents decolorization by acid-alcohol during staining (the acid-fast property), and inhibits lysosomal destruction within host phagocytes.

The transmission of \textit{M. tuberculosis} occurs through the inhalation of infectious droplet nuclei. These droplets, measuring 1 to 5 micrometres in diameter, are aerosolized when an individual with active pulmonary TB coughs, sneezes, speaks, or sings. Due to their small size, droplet nuclei bypass the mucociliary escalators of the upper respiratory tract and settle in the terminal alveoli of the lungs.

The immunopathological response of the host dictates the progression and clinical outcome of the infection, proceeding through the following stages:
\begin{itemize}
    \item \textbf{Alveolar Macrophage Phagocytosis:} Once the bacilli reach the alveoli, they are engulfed by resident alveolar macrophages. Under normal conditions, macrophages destroy pathogens by fusing the phagosome with a lysosome. \textit{Mycobacterium tuberculosis} avoids this by producing virulence factors (such as lipoarabinomannan and cord factor) that inhibit phagosome-lysosome fusion and block acidification of the phagosome, allowing the bacteria to survive and replicate intracellularly.
    \item \textbf{Macrophage Lysis and Chemokine Release:} As the intracellular bacilli replicate, they eventually cause the macrophage to lyse, releasing bacteria to infect surrounding cells. The damaged cells release cytokines and chemokines that recruit circulating monocytes, neutrophils, and lymphocytes to the site of infection, forming an early inflammatory focus.
    \item \textbf{Ghon Focus and Ranke Complex:} The primary site of parenchymal lung inflammation is typically in the subpleural region of the middle or lower lung zones, termed the Ghon focus. Bacilli migrate via local lymphatics to the draining hilar lymph nodes, causing lymphadenitis. The combination of the Ghon focus and the ipsilateral involved lymph node is known as the Ghon complex. If this complex undergoes calcification, it is referred to as a Ranke complex, visible on chest radiographs.
    \item \textbf{Cell-Mediated Immunity and Granuloma Formation:} Approximately 2 to 8 weeks post-infection, a cell-mediated immune response is activated. Antigen-presenting cells present mycobacterial peptides to CD4+ T helper cells, which differentiate into Th1 cells. These cells secrete pro-inflammatory cytokines, including Interferon-gamma (IFN-$\gamma$) and Tumor Necrosis Factor-alpha (TNF-$\alpha$). These cytokines activate macrophages, enhancing their bactericidal capacity and transforming them into epithelioid histiocytes, which may fuse to form multinucleated Langhans giant cells.
    To contain the infection, the immune system walls off the bacilli within a granuloma (tubercle). The granuloma is composed of a central core of infected and activated macrophages, surrounded by a rim of CD4+ and CD8+ T lymphocytes, B lymphocytes, and fibroblasts. The center of the granuloma undergoes caseous (cheese-like) necrosis, creating an acidic, hypoxic, and nutrient-deprived microenvironment that halts mycobacterial replication but permits the survival of dormant bacilli.
    \item \textbf{Latent TB Infection (LTBI):} In approximately 90\% to 95\% of immunocompetent hosts, the cell-mediated immune response successfully contains the bacilli within the granulomas. The infection becomes clinically latent, and the individual remains asymptomatic and non-infectious.
    \item \textbf{Reactivation (Post-Primary) TB:} In approximately 5\% to 10\% of individuals with LTBI, the latent bacilli reactivate, often decades later, due to a decline in cell-mediated immunity (e.g., aging, immunosuppressive medications, HIV infection, or malnutrition). Reactivation typically occurs in the apical and posterior segments of the upper lobes of the lungs, where high oxygen tension favors mycobacterial replication. The granulomas liquefy and rupture into the airways, causing extensive tissue destruction, cavitation, and clinical disease.
\end{itemize}

Key risk factors that predispose an individual to the development of active tuberculosis include:
\begin{itemize}
    \item \textbf{HIV Infection:} The most significant risk factor, as the depletion of CD4+ T lymphocytes directly impairs the host's ability to maintain granuloma integrity.
    \item \textbf{Immunosuppressive Medications:} These include systemic corticosteroids, chemotherapy, and tumor necrosis factor (TNF) antagonists (e.g., infliximab), which disrupt the cytokine pathways essential for granuloma maintenance.
    \item \textbf{Diabetes Mellitus:} Impairs neutrophil function, chemotaxis, and macrophage phagocytosis.
    \item \textbf{Severe Malnutrition:} Leads to generalized cell-mediated immunodeficiency.
    \item \textbf{Substance Abuse:} Chronic alcoholism, tobacco smoking, and intravenous drug use impair alveolar macrophage function and local pulmonary immunity.
    \item \textbf{Silicosis:} Inhaled silica particles damage alveolar macrophages, reducing their ability to clear mycobacteria.
    \item \textbf{Chronic Kidney Disease and End-Stage Renal Disease:} Uraemic toxins impair T-lymphocyte function.
\end{itemize}

\section*{Clinical Features}
The clinical presentation of tuberculosis is diverse and depends on the site of infection (pulmonary vs. extrapulmonary), the patient's immune status, and the extent of tissue involvement. The onset of symptoms is characteristically insidious, developing over weeks or months, which frequently delays medical evaluation.

Pulmonary tuberculosis represents 70\% to 80\% of active cases. The classic signs and symptoms include:
\begin{itemize}
    \item \textbf{Chronic Cough:} The most common clinical sign, typically presenting as a dry cough that becomes productive of mucoid, purulent, or blood-streaked sputum over a period of three weeks or more.
    \textbf{Haemoptysis:} The coughing up of blood, ranging from minor blood-streaking to massive, life-threatening airway haemorrhage. Massive haemoptysis is caused by the erosion of bronchial arteries or the rupture of a Rasmussen aneurysm (a dilated artery within a tuberculous cavity).
    \item \textbf{Fever:} A persistent, low-grade fever that classically exhibits a diurnal variation, rising in the late afternoon or evening.
    \item \textbf{Night Sweats:} Profuse, drenching sweats that occur during sleep, requiring the patient to change their clothing or bed linens.
    \item \textbf{Weight Loss and Anorexia:} Progressive weight loss accompanied by a loss of appetite, driven by chronic inflammatory cytokine release (such as TNF-$\alpha$).
    \item \textbf{Fatigue and Malaise:} Generalized physical exhaustion and a feeling of being unwell.
    \item \textbf{Chest Pain:} Pleuritic chest pain, which is worsened by inspiration or coughing, indicating pleural irritation or effusion.
\end{itemize}

Extrapulmonary tuberculosis occurs when the bacilli disseminate to other organ systems. This is more common in children and immunocompromised patients. Symptoms are specific to the involved organ system:
\begin{itemize}
    \item \textbf{Tuberculous Lymphadenitis (Scrofula):} Presents as painless, slowly enlarging, firm swelling of the lymph nodes, most commonly in the cervical region. The nodes may become matted, fluctuant, and eventually form chronic draining sinus tracts that discharge caseous material through the skin (scrofuloderma).
    \item \textbf{Pleural Tuberculosis:} Presents as an acute or subacute illness characterized by unilateral pleuritic chest pain, dry cough, and dyspnoea, caused by a pleural effusion secondary to hypersensitivity to mycobacterial antigens.
    \item \textbf{Skeletal Tuberculosis (Pott's Disease):} Tuberculous spondylitis typically involves the lower thoracic and upper lumbar spine. It causes progressive back pain, localized tenderness, stiffness, and vertebral destruction. This can lead to kyphosis (gibbus deformity) and spinal cord compression, resulting in sensory and motor deficits or paraplegia.
    \textbf{Central Nervous System Tuberculosis:} Includes tuberculous meningitis and tuberculomas. Tuberculous meningitis is a life-threatening emergency presenting with headache, low-grade fever, neck stiffness, altered mental status, confusion, and cranial nerve palsies (commonly III, VI, and VII).
    \item \textbf{Gastrointestinal and Peritoneal Tuberculosis:} Often involves the ileocecal region, presenting with abdominal pain, bloating, chronic diarrhoea, malabsorption, a palpable right lower quadrant mass, and ascites (peritoneal TB).
    \item \textbf{Genitourinary Tuberculosis:} Causes sterile pyuria (white blood cells in the urine with negative standard bacterial cultures), dysuria, haematuria, flank pain, and potential renal parenchymal destruction or ureteral strictures.
    \item \textbf{Miliary (Disseminated) Tuberculosis:} A life-threatening condition resulting from massive haematogenous dissemination of the bacilli. It presents with high, swinging fevers, rapid weight loss, hepatosplenomegaly, and respiratory distress. Chest radiographs show a classic pattern of tiny, bilateral nodules distributed uniformly throughout the lung fields, resembling millet seeds.
\end{itemize}

\section*{Differential Diagnosis}
Given its non-specific symptoms and variable presentation, tuberculosis can mimic many other pulmonary and systemic diseases. A comprehensive diagnostic workup must differentiate active pulmonary TB from the following conditions:
\begin{itemize}
    \item \textbf{Bacterial Pneumonia:} Acute community-acquired pneumonia (CAP) presents with fever, productive cough, and chest pain, but typically has a sudden onset (over days rather than weeks) and responds to standard empirical antibiotics.
    \item \textbf{Nontuberculous Mycobacterial (NTM) Lung Disease:} Caused by environmental mycobacteria (e.g., \textit{Mycobacterium avium} complex, \textit{Mycobacterium kansasii}). NTM lung disease clinically and radiographically mimics pulmonary TB, particularly in patients with pre-existing lung diseases such as COPD or bronchiectasis, but requires different antibiotic regimens.
    \item \textbf{Systemic Fungal Infections:} Fungal diseases such as Histoplasmosis, Coccidioidomycosis, and Blastomycosis can present with chronic cough, fever, weight loss, and cavitary lung lesions. These are typically associated with specific geographic exposures.
    \item \textbf{Bronchogenic Carcinoma (Lung Cancer):} Presents with chronic cough, haemoptysis, weight loss, and fatigue in older patients or smokers. Radiographs may reveal pulmonary masses, lobar collapse, or hilar lymphadenopathy, requiring cytological or histological exclusion of malignancy.
    \item \textbf{Sarcoidosis:} A systemic granulomatous disease characterized by non-caseating granulomas. It often presents with bilateral hilar lymphadenopathy, cough, and dyspnoea. Unlike TB, it does not show caseous necrosis on histopathology and does not respond to antitubercular therapy.
    \item \textbf{Pneumoconiosis (e.g., Silicosis):} Occupational lung disease caused by silica dust inhalation. It presents with nodular lung lesions on imaging. Silicosis impairs macrophage function, increasing the lifetime risk of developing tuberculosis (silicotuberculosis).
    \item \textbf{Bronchiectasis:} A chronic disease marked by permanent dilation of the bronchi, presenting with chronic productive cough, copious purulent sputum, and recurrent haemoptysis. It can mimic TB or occur as a long-term complication of past tuberculosis.
    \item \textbf{Lung Abscess:} A localized collection of pus in the lung parenchyma, typically caused by aspiration of anaerobic oral flora. It presents with high fever, cough with foul-smelling sputum, and a cavitary lesion on imaging, which can be mistaken for a tuberculous cavity.
    \item \textbf{Lymphoma:} Hodgkin and non-Hodgkin lymphomas can present with painless lymphadenopathy, night sweats, weight loss, and fever (Pel-Ebstein fever). Excisional lymph node biopsy is required to distinguish lymphoma from tuberculous lymphadenitis.
\end{itemize}

\section*{When to Seek Medical Care}
Active tuberculosis is a progressive disease that requires timely medical diagnosis and treatment. Early intervention is essential to prevent irreversible organ damage, improve clinical outcomes, and halt the transmission of the airborne bacilli to family members, co-workers, and the community.

A general practitioner or primary care doctor should be contacted for an appointment if an individual experiences any of the following persistent, unexplained symptoms:
\begin{itemize}
    \item A cough that lasts for three weeks or longer.
    \item Unexplained, progressive weight loss and loss of appetite.
    \item Drenching night sweats.
    \item A persistent, low-grade fever, particularly in the late afternoon or evening.
    \item Generalized, unexplained fatigue and lack of energy.
    \item Painless swelling of lymph nodes, particularly in the neck or supraclavicular region.
\end{itemize}

Certain severe, acute symptoms indicate a medical emergency, representing high-risk complications of tuberculosis or rapid disease progression. Immediate emergency medical services must be contacted (for example, by calling \textbf{local emergency services} in many countries, \textbf{911} in the United States and Canada, or \textbf{999} in the United Kingdom) if a patient exhibits:
\begin{itemize}
    \item \textbf{Severe Difficulty Breathing:} Shortness of breath, rapid respiration, or gasping for air, which may indicate a massive pleural effusion, extensive lung destruction, or acute respiratory distress syndrome (ARDS).
    \item \textbf{Massive Haemoptysis:} Coughing up large volumes of bright red blood (typically defined as more than 100 to 200 mL in a 24-hour period), which carries a high risk of asphyxiation and hypovolaemic shock.
    \item \textbf{Signs of Meningitis:} A sudden, severe headache accompanied by a stiff neck (inability to touch the chin to the chest), high fever, confusion, photophobia (sensitivity to light), drowsiness, or a change in mental state.
    \item \textbf{Sudden Neurological Deficits:} Weakness, numbness, loss of sensation, or paralysis in the limbs, which may indicate spinal cord compression from skeletal TB (Pott's disease) or central nervous system lesions.
    \item \textbf{Chest Pain:} Sudden, sharp, or severe chest pain that is exacerbated by breathing or coughing, which can indicate pleural inflammation, pneumothorax, or pericardial involvement.
\end{itemize}

\section*{Diagnosis and Investigations}
The diagnosis of tuberculosis requires a multi-faceted approach combining clinical evaluation, immunological screening for exposure, radiological imaging, and microbiological confirmation of the causative pathogen. Because of the public health implications, establishing a definitive diagnosis is paramount.

\subsection*{Physical Examination}
The physical examination is often non-specific. In pulmonary TB, auscultation of the chest may reveal decreased breath sounds, crackles (crepitations), or bronchial breathing over areas of consolidation or cavitation. Cachexia, fever, and pallor (secondary to anemia of chronic disease) may be observed. Examination of the neck, axilla, and inguinal regions is crucial to check for lymphadenopathy, while neurological examinations are necessary to rule out signs of meningeal irritation or focal deficits.

\subsection*{Immunological Screening for Latent TB Infection (LTBI)}
These tests detect the host's immune response to mycobacterial antigens. They cannot distinguish between latent infection and active disease:
\begin{itemize}
    \item \textbf{Tuberculin Skin Test (TST) / Mantoux Test:} Involves an intradermal injection of 5 tuberculin units (0.1 mL) of Purified Protein Derivative (PPD) into the volar aspect of the forearm. The test is read 48 to 72 hours later. The diameter of the induration (not erythema) is measured in millimetres. The threshold for a positive result depends on the patient's risk profile: $\ge 5$ mm for high-risk patients (e.g., HIV-infected, close contacts of active cases), $\ge 10$ mm for moderate-risk patients (e.g., immigrants from high-burden areas, healthcare workers), and $\ge 15$ mm for low-risk individuals. A major limitation of TST is cross-reactivity with the BCG vaccine and nontuberculous mycobacteria, leading to false-positive results.
    \item \textbf{Interferon-Gamma Release Assays (IGRAs):} In vitro blood tests (such as QuantiFERON-TB Gold Plus and T-SPOT.TB) that measure the amount of IFN-$\gamma$ released by T lymphocytes when mixed with specific \textit{M. tuberculosis} antigens (ESAT-6 and CFP-10). These antigens are absent from the BCG vaccine strain and most NTMs, making IGRAs highly specific and unaffected by prior BCG vaccination. IGRAs require only a single patient visit but are more costly than TST.
\end{itemize}

\subsection*{Radiological Imaging}
\begin{itemize}
    \item \textbf{Chest Radiography (CXR):} The primary imaging modality for suspected pulmonary TB. In primary progressive TB, CXR typically shows lower-lobe consolidations, hilar lymphadenopathy, or pleural effusions. In reactivation (post-primary) TB, the classic findings include patchy infiltrates or consolidations in the apical and posterior segments of the upper lobes or the superior segment of the lower lobes. Cavitation (lucent areas with thick walls) is a hallmark of highly infectious disease. Fibrotic scarring and calcification indicate past, healed infection.
    \item \textbf{Computed Tomography (CT) of the Chest:} More sensitive than CXR for detecting early infiltrates, small cavities, mediastinal lymphadenopathy, and the ``tree-in-bud'' pattern, which signifies active bronchogenic spread of infection.
\end{itemize}

\subsection*{Microbiological Confirmation (The Gold Standard)}
A definitive diagnosis of active TB requires the demonstration of \textit{M. tuberculosis} in clinical specimens (sputum, induced sputum, bronchoalveolar lavage fluid, pleural fluid, cerebrospinal fluid, urine, or tissue biopsy):
\begin{itemize}
    \item \textbf{Acid-Fast Bacilli (AFB) Smear Microscopy:} Clinical specimens are stained using the Ziehl-Neelsen method or auramine-rhodamine fluorescent staining. Smeared slides are examined under a microscope to detect acid-fast bacilli. Although rapid (results within hours) and inexpensive, microscopy has low sensitivity (requiring at least 5,000 to 10,000 bacilli per mL of sputum to be positive) and cannot differentiate \textit{M. tuberculosis} from NTMs.
    \item \textbf{Mycobacterial Culture:} Sputum is inoculated onto solid media (Lowenstein-Jensen slants) or liquid media (Mycobacteria Growth Indicator Tube - MGIT). Culturing is the gold standard because it is highly sensitive (detecting as few as 10 to 100 bacilli per mL) and allows for subsequent Drug Susceptibility Testing (DST). However, due to the slow growth rate of the bacterium, solid culture can take 3 to 8 weeks, while liquid culture takes 10 to 21 days.
    \item \textbf{Nucleic Acid Amplification Tests (NAATs):} Molecular assays such as the GeneXpert MTB/RIF test. These tests amplify and detect \textit{M. tuberculosis} DNA directly from sputum samples within 2 hours. GeneXpert also identifies mutations in the \textit{rpoB} gene, which mediates resistance to rifampicin, providing a rapid screen for drug-resistant strains. The WHO recommends GeneXpert as the initial diagnostic test for all individuals suspected of having active TB.
\end{itemize}

\subsection*{Histopathological Examination}
Biopsy specimens of lymph nodes, pleura, or other affected tissues show characteristic caseating granulomas. The presence of focal central necrosis (caseation) surrounded by epithelioid histiocytes, Langhans giant cells, and a peripheral lymphocyte rim strongly suggests tuberculosis, even if AFB staining of the tissue is negative.

\section*{Management}
The management of tuberculosis involves a structured, prolonged course of combination antimicrobial therapy designed to eradicate the replicating bacilli, prevent the development of drug resistance, and sterilize dormant populations to prevent relapse. Treatment must be coordinated by public health authorities and experienced clinicians.

\subsection*{Pharmacotherapy for Drug-Susceptible Active Pulmonary TB}
The standard regimen for drug-susceptible tuberculosis spans 6 months and is divided into two distinct phases:
\begin{itemize}
    \item \textbf{Intensive Phase (2 Months):} Designed to rapidly kill the actively replicating bacilli, reduce the bacterial load, improve symptoms, and halt transmission. The regimen consists of four first-line drugs administered daily:
    \begin{itemize}
        \item \textbf{Isoniazid (H):} Highly bactericidal, targeting cell-wall mycolic acid synthesis.
        \item \textbf{Rifampicin (R):} Bactericidal, inhibiting bacterial RNA synthesis (RNA polymerase).
        \item \textbf{Pyrazinamide (Z):} Sterilizing agent active in the acidic environment of granulomas.
        \item \textbf{Ethambutol (E):} Bacteriostatic, inhibiting arabinosyltransferase involved in cell-wall synthesis, added primarily to prevent the emergence of resistance if isoniazid resistance is present.
    \end{itemize}
    This phase is commonly abbreviated as the 2-month HRZE regimen.
    \item \textbf{Continuation Phase (4 Months):} Designed to eliminate the semi-dormant ``persister'' bacilli within granulomas, ensuring complete sterilization and preventing subsequent relapse. The regimen consists of two drugs:
    \begin{itemize}
        \item \textbf{Isoniazid (H)}
        \item \textbf{Rifampicin (R)}
    \end{itemize}
    This phase is abbreviated as the 4-month HR regimen. The total course is written as 2HRZE/4HR.
\end{itemize}

The continuation phase is extended from 4 to 7 months (making the total treatment 9 months) in specific scenarios: if the patient has cavitary pulmonary TB at baseline and remains sputum culture-positive at the end of the 2-month intensive phase; if pyrazinamide was not included in the intensive phase; or in cases of central nervous system (CNS) tuberculosis or osteoarticular tuberculosis.

\subsection*{Directly Observed Therapy (DOTS)}
Non-adherence to TB therapy is the primary driver of treatment failure, relapse, and the development of drug-resistant strains. To address this, the WHO advocates for Directly Observed Therapy (DOTS), a strategy where a trained healthcare provider or designated community member observes the patient swallow every single dose of medication. Modern variations include Video-Observed Therapy (VOT), where patients record and upload videos of themselves taking their medications via smartphone applications.

\subsection*{Monitoring and Managing Drug Toxicity}
Due to the hepatotoxic nature of isoniazid, rifampicin, and pyrazinamide, patients must undergo baseline and periodic Liver Function Tests (LFTs), monitoring serum transaminases (AST, ALT) and bilirubin.
\begin{itemize}
    \item \textbf{Hepatotoxicity:} If transaminase levels increase to $>3$ times the upper limit of normal in the presence of symptoms (such as nausea, vomiting, abdominal pain, or jaundice), or $>5$ times the upper limit of normal in asymptomatic patients, all hepatotoxic drugs must be suspended immediately. Medications are sequentially reintroduced once liver enzymes normalize.
    \item \textbf{Peripheral Neuropathy:} Isoniazid interferes with pyridoxine (Vitamin B6) metabolism, which can lead to peripheral neuropathy (characterized by numbness, tingling, or burning sensations in the hands and feet). To prevent this, Pyridoxine supplementation (10 to 25 mg daily) is routinely co-administered to patients at high risk (e.g., those with diabetes, malnutrition, alcoholism, chronic kidney disease, pregnancy, or HIV).
    \item \textbf{Optic Neuritis:} Ethambutol can cause dose-dependent, reversible or irreversible optic neuritis, presenting with decreased visual acuity and loss of red-green color discrimination. Patients should undergo baseline and monthly assessments of visual acuity and color vision.
    \item \textbf{Hyperuricaemia:} Pyrazinamide inhibits renal excretion of uric acid, frequently causing asymptomatic hyperuricaemia. While mild arthralgias are common and managed with NSAIDs, acute gouty arthritis requires cessation of the drug.
    \item \textbf{Drug Interactions:} Rifampicin is a potent inducer of hepatic cytochrome P450 enzymes (specifically CYP3A4). It accelerates the metabolism of numerous medications, rendering them less effective. Important interactions include oral contraceptives, warfarin, anticonvulsants, antiretroviral drugs (necessitating adjustments in HIV regimens), and oral anticoagulants.
\end{itemize}

\subsection*{Treatment of Latent TB Infection (LTBI)}
Treating LTBI is crucial to prevent the progression to active disease. Several safe and effective regimens are available:
\begin{itemize}
    \item Isoniazid daily for 6 or 9 months (6H or 9H).
    \item Rifampicin daily for 4 months (4R).
    \item Isoniazid and Rifampicin daily for 3 months (3HR).
    \item Weekly Isoniazid and Rifapentine for 3 months (3HP), administered under direct observation.
\end{itemize}

\subsection*{Treatment of Drug-Resistant TB}
MDR-TB and XDR-TB require complex, individualized regimens utilizing second-line antitubercular drugs. These regimens are longer (ranging from 9 to 20 months), less effective, and associated with higher rates of toxicity. Regimens incorporate drugs such as fluoroquinolones (levofloxacin, moxifloxacin), bedaquiline, linezolid, pretomanid, cycloserine, and clofazimine. Injectable agents (amikacin, kanamycin) have largely been phased out due to irreversible ototoxicity and nephrotoxicity.

\section*{Complications}
Active tuberculosis can cause severe, irreversible tissue destruction, leading to systemic and organ-specific complications, particularly if diagnosis and treatment are delayed.

\begin{itemize}
    \item \textbf{Permanent Pulmonary Damage:} Extensive parenchymal destruction, necrosis, and subsequent healing by fibrosis can leave patients with permanent lung damage. This includes bronchiectasis (destruction and widening of the large airways), parenchymal scarring, and chronic pleural thickening (fibrothorax). These structural changes result in restrictive or obstructive lung disease, presenting as chronic dyspnoea, recurrent secondary bacterial infections, and chronic respiratory failure.
    \item \textbf{Haemoptysis and Rasmussen Aneurysms:} Persistent cavitary lesions can erode into pulmonary blood vessels. A Rasmussen aneurysm represents a localized dilation of a branch of the pulmonary artery within a tuberculous cavity. Rupture of this aneurysm causes massive, catastrophic haemoptysis, which is a life-threatening airway emergency.
    \item \textbf{Pneumothorax and Empyema:} A tuberculous cavity situated near the periphery of the lung can rupture into the pleural space. This allows air to enter the pleural cavity, causing a pneumothorax (lung collapse), or introduces caseous, purulent material, resulting in a tuberculous empyema (pus in the pleural space) or bronchopleural fistula.
    \item \textbf{Cor Pulmonale:} Chronic, extensive destruction of the pulmonary vascular bed due to widespread fibroid scarring leads to pulmonary hypertension. Over time, this increases the workload on the right side of the heart, resulting in right ventricular hypertrophy and eventually right-sided heart failure (cor pulmonale).
    \item \textbf{Neurological Deficits:} In tuberculous meningitis, the inflammatory exudate at the base of the brain can entrap cranial nerves, leading to palsies (commonly affecting nerves III, VI, and VII). Obstruction of the cerebrospinal fluid (CSF) flow path leads to communicating or non-communicating hydrocephalus. Ischaemic strokes can occur due to vasculitis of the cerebral blood vessels, resulting in permanent cognitive impairment, hemiparesis, or aphasia.
    \item \textbf{Spinal Deformity and Paraplegia:} Pott's disease (skeletal TB) causes progressive destruction of the intervertebral discs and adjacent vertebral bodies. This leads to collapse of the anterior spine, resulting in kyphosis (gibbus deformity). The collapse can compress the spinal cord, causing sensory loss, bowel/bladder dysfunction, and paraplegia.
    \item \textbf{Adrenal Insufficiency (Addison's Disease):} Tuberculous destruction of the adrenal glands was historically the leading cause of Addison's disease and remains a significant cause in endemic countries. It presents with chronic fatigue, hyperpigmentation of the skin, hypotension, electrolyte imbalances, and life-threatening adrenal crisis.
    \item \textbf{Amyloidosis:} Chronic, uncontrolled systemic inflammation in untreated tuberculosis can lead to secondary (AA) amyloidosis. This condition involves the extracellular deposition of amyloid fibril proteins in organs such as the kidneys, liver, and spleen, leading to nephrotic syndrome and progressive renal failure.
\end{itemize}

\section*{Prognosis and Outcomes}
The prognosis for individuals diagnosed with tuberculosis is highly dependent on several factors: the drug susceptibility of the mycobacterial strain, the timing of treatment initiation, patient adherence to the prescribed medication regimen, and the presence of underlying comorbidities (especially HIV infection).

For drug-susceptible tuberculosis, the prognosis is excellent if a full, uninterrupted course of combination therapy is completed. Cure rates exceeding 95\% are routinely achieved in clinical trials and well-structured public health programs. Clinical improvement is typically rapid; fever and constitutional symptoms resolve within 2 to 4 weeks of starting therapy, and sputum conversion (smears and cultures becoming negative for \textit{M. tuberculosis}) occurs in over 80\% of patients by the end of the 2-month intensive phase.

Conversely, untreated active tuberculosis carries a grave prognosis. Historically, before the advent of chemotherapy, approximately 50\% of patients with active pulmonary TB died within five years of diagnosis, and only 25\% achieved spontaneous recovery, while the remaining 25\% remained chronically ill, serving as continuous sources of transmission.

The prognosis is significantly worsened in the following cohorts:
\begin{itemize}
    \item \textbf{HIV Co-infection:} Immunocompromised patients have a much higher risk of rapid disease progression, disseminated (miliary) disease, and death. Without antiretroviral therapy and concurrent TB treatment, mortality is extremely high. However, integrated care addressing both infections dramatically improves survival.
    \item \textbf{Drug-Resistant TB:} Treatment success rates for MDR-TB are lower, hovering around 60\% globally, while success rates for XDR-TB are lower still (approximately 30\% to 40\%). These patients face prolonged courses of toxic medications, and failure to cure can lead to chronic, drug-resistant infection and high mortality.
    \item \textbf{Extremes of Age:} Infants, young children, and the frail elderly are at higher risk of severe forms of TB (such as meningitis and miliary TB) and have higher mortality rates.
\end{itemize}
Following successful completion of therapy, patients require monitoring for relapse, which occurs in less than 5\% of patients treated with the standard 6-month regimen, typically within the first two years post-treatment. Residual lung damage may persist, requiring long-term supportive pulmonary care.

\section*{Prevention and Self-Care}
Controlling tuberculosis requires a combination of public health strategies, clinical interventions, and individual preventative measures. Because the disease is airborne, prevention must focus on reducing transmission risks alongside boosting host immunity.

\subsection*{Public Health and Medical Interventions}
\begin{itemize}
    \item \textbf{Vaccination (BCG Vaccine):} The Bacille Calmette-Gu\'{e}rin (BCG) vaccine is a live-attenuated strain derived from \textit{Mycobacterium bovis}. It is the only licensed vaccine for tuberculosis. Administered intradermally at birth, the BCG vaccine provides robust protection (up to 80\%) against severe, disseminated forms of TB in infants and young children, such as tuberculous meningitis and miliary TB. However, its efficacy in preventing adult pulmonary tuberculosis is highly variable (ranging from 0\% to 80\%) and tends to wane over time. Consequently, it is routinely given in high-burden countries but is only recommended for select high-risk populations in low-burden countries.
    \item \textbf{Infection Control in Healthcare Settings:} Hospitals and clinics managing TB patients must employ strict engineering and administrative controls. This includes placing infectious patients in negative-pressure isolation rooms, ensuring adequate air exchanges (natural or mechanical ventilation), and utilizing ultraviolet germicidal irradiation (UVGI) to kill airborne bacilli. Healthcare workers must wear fit-tested N95 or particulate respirators when entering these zones.
    \item \textbf{Contact Tracing and Screening:} Public health departments perform contact tracing whenever an active case of pulmonary TB is identified. Close contacts (family, housemates, co-workers) are screened using TST or IGRAs, alongside chest X-rays. Those diagnosed with active TB are immediately treated, while those with LTBI are offered preventative therapy to abort progression.
\end{itemize}

\subsection*{Self-Care and Personal Prevention}
Patients undergoing treatment and individuals seeking to protect themselves should adopt the following self-care measures:
\begin{itemize}
    \item \textbf{Strict Treatment Adherence:} The single most important action for a patient diagnosed with TB is to take their medications exactly as prescribed. Doses must not be missed, and the course must not be shortened, even if symptoms resolve completely. Erratic drug intake leads to the selection of drug-resistant mutants, which are far more difficult and dangerous to treat.
    \item \textbf{Respiratory Hygiene:} While infectious (usually the first 2 to 3 weeks of treatment for drug-susceptible TB), patients should stay home from work, school, and public places. They should sleep in a separate, well-ventilated room. When coughing or sneezing, the mouth and nose must be covered with a tissue, which should be immediately discarded in a closed bin, or into the elbow.
    \item \textbf{Optimizing Ventilation:} Keeping windows open to allow fresh air to circulate reduces the concentration of droplet nuclei in indoor environments, drastically reducing transmission risks.
    \item \textbf{Nutritional Support:} TB causes significant wasting. Eating a balanced diet rich in protein, calories, vitamins, and minerals supports immune function and aids in tissue recovery. Nutritional counseling and support are integral parts of comprehensive TB care.
    \item \textbf{Avoiding Alcohol and Tobacco:} Alcohol consumption increases the risk of drug-induced hepatotoxicity from TB medications and must be avoided. Tobacco smoke damages the respiratory epithelium, impairs alveolar macrophage function, worsens symptoms, and increases the risk of treatment failure and relapse.
    \item \textbf{Monitoring for Adverse Effects:} Patients must be educated to recognize the signs of drug toxicity (such as yellowing of the skin/eyes, persistent nausea, dark urine, joint pain, or visual changes) and contact their healthcare provider immediately if they occur.
\end{itemize}

\section*{Sources and Further Reading}
For authoritative, evidence-based information and updates on tuberculosis guidelines, resources from the following international and national organizations are recommended:
\begin{itemize}
    \item World Health Organization (WHO) Global Tuberculosis Programme: \url{https://www.who.int/teams/global-tuberculosis-programme}
    \item Centers for Disease Control and Prevention (CDC) Division of Tuberculosis Elimination: \url{https://www.cdc.gov/tb/index.html}
    \item National Institute for Health and Care Excellence (NICE) Tuberculosis Guideline [NG33]: \url{https://www.nice.org.uk/guidance/ng33}
    \item International Union Against Tuberculosis and Lung Disease (The Union): \url{https://theunion.org/}
    \item Stop TB Partnership: \url{https://www.stoptb.org/}
    \item Australian Government Department of Health - National Tuberculosis Advisory Committee (NTAC): \url{https://www.health.gov.uk/committees-and-groups/national-tuberculosis-advisory-committee-ntac}
\end{itemize}